%! Author = sriadityavedantam
%! Date = 3/30/26

\section{Normal Subgroups}

We have previously defined left congruence as follows.
Let $H\leq G$ and $a,b\in G$.
Then
\[
    a\overset{\ell}{\equiv} b\mod H
\]
if
\[
    \vs{a}b\in H.
\]
We also defined left cosets as
\[
    aH\coloneqq \{ah:h\in H\}.
\]

\begin{definition}[Right Congruence]
    Let $H\leq G$.
    Define
    \[
        a\overset{r}{\equiv} b\mod H
    \]
    if
    \[
        b\vs{a}\in H.
    \]
    Equivalently,
    \[
        a\vs{b}\in H.
    \]
    This defines an equivalence relation.
    The equivalence classes are called right cosets and are given by
    \[
        Ha\coloneqq \{ha:h\in H\}.
    \]
\end{definition}

Note that
\[
    a\in aH\cap Ha,
\]
but this may be the only common element.
In general,
\[
    aH\neq Ha.
\]

\begin{example}
    Let
    \[
        G=S_3,
        \quad
        H=\{e,(12)\}=\gen{(12)}.
    \]
\end{example}

Left cosets:
\[
    H=\{e,(12)\},
    \quad
    (13)H=\{(13),(132)\},
    \quad
    (23)H=\{(23),(123)\}.
\]
Right cosets:
\[
    H=\{e,(12)\},
    \quad
    H(13)=\{(13),(123)\},
    \quad
    H(23)=\{(23),(132)\}.
\]

\begin{example}
    Let
    \[
        N=\{e,(123),(132)\}=\gen{(123)}.
    \]
\end{example}

Left cosets:
\[
    N,
    \quad
    (12)N=\{(12),(23),(13)\}=(23)N=(13)N.
\]
Right cosets:
\[
    N,
    \quad
    N(12).
\]
Thus
\[
    [G:N]=2.
\]
Recall that the number of left cosets always equals the number of right cosets:
\[
    [G:H]=\frac{\ord(G)}{\ord(H)}.
\]

\begin{definition}[Normal Subgroup]
    Let $N\leq G$.
    We say $N$ is normal in $G$ if the left cosets equal the right cosets.
    This is denoted by
    \[
        N\trianglelefteq G.
    \]
    That is,
    \[
        aN=Na
        \quad\text{for all }a\in G.
    \]
\end{definition}

\begin{theorem}{
    Let $N\leq G$.
    The following are equivalent:
    \begin{enumerate}
        \item The set of left cosets equals the set of right cosets.
        \item For all $a\in G$, $aN=Na$.
        \item For all $a\in G$, $aN\vs{a}=N$.
        \item For all $a\in G$, $\vs{a}Na=N$.
        \item For all $a\in G$ and $x\in N$, $axa^{-1}\in N$.
        \item For all $a\in G$ and $x\in N$, $a^{-1}xa\in N$.
        \item Multiplication of left cosets is well-defined, i.e.
        \[
            (aN)(bN)=abN.
        \]
        \item Multiplication of right cosets is well-defined.
    \end{enumerate}
}
\end{theorem}

\begin{proof}
    We sketch key implications.

    (2)$\implies$(1) is immediate.

    (1)$\implies$(2):
    If left and right cosets coincide, then for any $a\in G$, the left coset $aN$ must equal the right coset containing $a$, so $aN=Na$.

    (2)$\implies$(5):
    Let $x\in N$.
    Then $ax\in aN=Na$, so $ax=ya$ for some $y\in N$.
    Thus
    \[
        axa^{-1}=y\in N.
    \]

    (5)$\implies$(4):
    For all $x\in N$, we have $axa^{-1}\in N$, so $aNa^{-1}\subseteq N$.
    Applying the same argument with $a^{-1}$ gives equality.

    (2)$\implies$(7):
    Suppose $aN=cN$ and $bN=dN$.
    Then $\vs{c}a\in N$ and $\vs{d}b\in N$.
    Using normality,
    \[
        \vs{d}\vs{c}ab\in N,
    \]
    so
    \[
        abN=cdN.
    \]
    Thus coset multiplication is well-defined.
\end{proof}

\begin{theorem}{
    If $N\trianglelefteq G$, then multiplication of cosets is well-defined.
    That is, if $aN=cN$ and $bN=dN$, then
    \[
        abN=cdN.
    \]
}
\end{theorem}

\begin{proof}
    If $aN=cN$, then $\vs{c}a\in N$.
    If $bN=dN$, then $\vs{d}b\in N$.
    Since $N$ is normal,
    \[
        \vs{d}(\vs{c}a)d\in N.
    \]
    Multiplying,
    \[
        (\vs{d}(\vs{c}a)d)(\vs{d}b)=\vs{d}\vs{c}ab\in N.
    \]
    Thus
    \[
        \vs{(cd)}ab\in N,
    \]
    so $abN=cdN$.
\end{proof}

\begin{proposition}{
    $N\trianglelefteq G$ if and only if for all $a\in G$ and $x\in N$,
    \[
        axa^{-1}\in N.
    \]
}
\end{proposition}

\begin{proof}
    $(\implies)$
    If $N$ is normal, then $aN=Na$.
    So for $x\in N$, $ax\in Na$, hence $ax=ya$ for some $y\in N$.
    Thus
    \[
        axa^{-1}=y\in N.
    \]
    $(\impliedby)$
    If $axa^{-1}\in N$ for all $a\in G$ and $x\in N$, then $aN\subseteq Na$.
    A similar argument gives $Na\subseteq aN$, so $aN=Na$.
\end{proof}

\begin{definition}[Characteristic Subgroup]
    A subgroup $N\leq G$ is characteristic in $G$ if for every automorphism $\phi\in\aut(G)$,
    \[
        \phi(N)=N.
    \]
    This is denoted by
    \[
        N\char G.
    \]
\end{definition}

\begin{lemma}{
    If $N\char G$, then $N\trianglelefteq G$.
}
\end{lemma}

\begin{lemma}{
    The center of $G$ is characteristic in $G$.
}
\end{lemma}

\begin{definition}[Quotient Group]
    Let $N\trianglelefteq G$.
    The quotient group is
    \[
        G/N\coloneqq \{gN:g\in G\}.
    \]
    The operation is defined by
    \[
        (gN)(hN)=ghN.
    \]
\end{definition}